\documentclass[11pt]{article}
\usepackage[margin=1in]{geometry}
\usepackage{amsmath,amssymb,amsthm}
\usepackage{booktabs}
\usepackage[round]{natbib}
\usepackage{hyperref}
\usepackage{microtype}

\newtheorem{theorem}{Theorem}
\newtheorem{proposition}{Proposition}
\newtheorem{lemma}{Lemma}
\newtheorem{corollary}{Corollary}
\theoremstyle{definition}
\newtheorem{example}{Example}
\theoremstyle{remark}
\newtheorem{remark}{Remark}

\newcommand{\E}{\mathbb{E}}
\newcommand{\Prob}{\mathbb{P}}
\newcommand{\Var}{\operatorname{Var}}
\newcommand{\set}[1]{\{#1\}}

\title{Conformal Correlation Counterexamples}
\author{Peter Cotton\\\texttt{peter.cotton@microprediction.com}}
\date{September 2026}

\begin{document}
\maketitle

\begin{abstract}
The conformal correlation matrix of \citet{perlo2026} is the phi coefficient between indicators of
class membership in conformal prediction sets. The authors read it as a map of which classes a
classifier confuses. Eight hand-checkable examples show that it is not. Two classifiers built inside
randomized adaptive prediction sets, exactly as specified, produce the same law of prediction sets
and the same matrix while confusing classes in opposite directions, so the matrix does not
determine the confusion matrix. The reason is that the matrix depends only on the law of the
prediction set, while confusion depends on the joint law of the set, the true label and the ranked
prediction. Two lemmas give every entry for classes that never share a set as a function of the
inclusion rates alone, and recover every inclusion rate exactly from three entries when sets are
singletons, which bears on the privacy claim made for the matrix.
\end{abstract}

\section{The statistic}\label{sec:statistic}

There are $H$ classes. A classifier maps an input $X$ to a probability vector $p(X)$. Split
conformal prediction \citep{vovk2005, lei2018} turns the probability vector into a prediction set
$C(X) \subseteq \set{1, \dots, H}$. The true label is $Y$ and the class the classifier ranks first is
$\hat Y$.

The authors build their sets with the adaptive prediction sets (APS) of \citet{romano2020}. APS
assigns a score to each label at each input. Sort the probabilities at $x$ in decreasing order,
$p_{(1)}(x) \ge \dots \ge p_{(H)}(x)$, write $r(x, y)$ for the rank of class $y$ in that order, and
write $S_k(x) = \sum_{l \le k} p_{(l)}(x)$ for the probability mass through rank $k$. The score of
label $y$ at input $x$ is
\begin{equation}\label{eq:aps}
s(x, y, u) = S_{r(x,y)}(x) - u\, p_y(x), \qquad u \sim \mathrm{Unif}(0, 1),
\end{equation}
the mass at and above $y$ in the ranking, less a uniform share of the mass of $y$ itself.

The score is turned into a set by calibration. Given $m$ labeled calibration pairs $(x_i, y_i)$
with independent uniforms $u_i$, let $\tau$ be the $\lceil (m+1)(1-\alpha) \rceil$-th smallest of
the calibration scores $s(x_i, y_i, u_i)$. The prediction set at a new input $x$, with a fresh
uniform $U$, is
\begin{equation}\label{eq:apsset}
C(x) = \set{\, y : s(x, y, U) \le \tau \,}.
\end{equation}
In words, the set is the top ranked classes down to the first rank at which the cumulative mass
reaches $\tau$, and the class at that boundary rank is kept or dropped according to $U$.

Two coverage bounds hold. Under exchangeability of the calibration pairs and the test pair,
$\Prob(Y \in C(X)) \ge 1 - \alpha$. When the scores are almost surely distinct, also
$\Prob(Y \in C(X)) \le 1 - \alpha + 1/(m+1)$ \citep[Theorem~1]{romano2020}. Setting $u = 0$
throughout gives the non randomized variant, which is conservative. Example~\ref{ex:aps} uses the
randomized variant exactly as defined here. Nothing else below depends on APS, beyond the fact that
$C(X)$ is a set of classes.

The matrix is built from membership indicators. For each class $h$ define
\[
z_h = \mathbf{1}\set{h \in C(X)},
\]
and write $p_h = \E[z_h]$ for the inclusion rate of class $h$ and $q_{ij} = \E[z_i z_j]$ for the
co-occurrence rate of classes $i$ and $j$. The conformal correlation matrix (CCM) of
\citet{perlo2026} has entries
\begin{equation}\label{eq:ccm}
\rho_{ij} = \frac{\E[z_i z_j] - \E[z_i]\,\E[z_j]}{\sqrt{\Var(z_i)\,\Var(z_j)}}
          = \frac{q_{ij} - p_i p_j}{\sqrt{p_i(1-p_i)\,p_j(1-p_j)}} .
\end{equation}
This is the Pearson correlation of two Bernoulli variables, also called the phi coefficient
\citep{yule1912}. In practice the expectations are sample averages over a test set.

The entry is defined only when $0 < p_i < 1$ and $0 < p_j < 1$, since otherwise a variance is zero.
The authors give no convention for that case. We exclude such classes from the matrix.

The authors attach three readings to the entries. Their Remark~1 says that intrinsically similar
classes ``might be expected to exhibit high positive correlation''. Their Remark~2 says that pairs of
dissimilar classes ``are expected to display strong negative correlation''. Their Remark~3 says that
classes ``unrelated or independent with respect to the classifier representation are expected to
yield correlation coefficients close to zero''.

On the strength of these readings they draw three conclusions. A positive entry between two classes
means the classifier ``is much more likely to misclassify'' one as the other. The matrix ``contains a
superset of the information conveyed by the confusion matrix''. And the matrix ``captures the
likelihood of class pairs co-occurring within conformal prediction sets''. Each of the three is
false, and the examples below show it.

\section{The matrix depends only on the law of the sets}\label{sec:organizing}

\begin{proposition}\label{prop:law}
The conformal correlation matrix is a functional of the law of the prediction set $C(X)$ alone.
Two situations with the same law of $C(X)$ have the same matrix, whatever the joint law of $(C(X),
Y, \hat Y)$.
\end{proposition}

\begin{proof}
Every quantity in \eqref{eq:ccm} is an expectation of a function of the vector $(z_1, \dots, z_H)$,
and that vector is a function of $C(X)$.
\end{proof}

Confusion is a different object. The confusion matrix has entries $\Prob(Y = a, \hat Y = b)$, which
depend on the joint law of the true label and the ranked prediction. Hold the law of $C(X)$ fixed
and change $Y$ or $\hat Y$, and by Proposition~\ref{prop:law} the matrix cannot notice.

The only thing that ties $C(X)$ to $Y$ is coverage. Under exchangeability of calibration and test
data, split conformal prediction guarantees
\begin{equation}\label{eq:coverage}
\Prob(z_Y = 1) = \sum_{y} \Prob(Y = y,\; z_y = 1) \;\ge\; 1 - \alpha .
\end{equation}
This is one number. It constrains the sum on the right and nothing else.

In particular, coverage does not in general determine a co-occurrence rate $q_{ij}$. That rate mixes
the true class with a wrong class, two wrong classes, and different true classes on different
inputs, and coverage never justifies reading it as confusion. There is one special case worth
stating. With only two classes and coverage one the set is never empty, so
$0 = \Prob(C = \varnothing) = 1 - p_A - p_B + q_{AB}$ and $q_{AB} = p_A + p_B - 1$ is forced. With
three or more classes nothing of the kind holds, as the next proposition shows for a single entry.

\begin{proposition}\label{prop:any}
Among laws of prediction sets that satisfy the marginal coverage lower bound \eqref{eq:coverage},
for every $\rho \in [-1, 1]$ there is one with coverage one and accuracy one whose matrix has an
off-diagonal entry equal to $\rho$.
\end{proposition}

\begin{proof}
Let the true label always be $T$, always in the set and always ranked first, so coverage and
accuracy are one. Since $z_T \equiv 1$ the entries involving $T$ are undefined, and the claim
concerns $\rho_{AB}$. Let two other labels $A$ and $B$ be added to the set according to
\begin{align*}
\Prob(z_A = 0, z_B = 0) &= \Prob(z_A = 1, z_B = 1) = \frac{1+\rho}{4}, \\
\Prob(z_A = 1, z_B = 0) &= \Prob(z_A = 0, z_B = 1) = \frac{1-\rho}{4}.
\end{align*}
Then $p_A = p_B = \tfrac{1+\rho}{4} + \tfrac{1-\rho}{4} = \tfrac12$, so each variance is
$\tfrac14$, and $q_{AB} = \tfrac{1+\rho}{4}$. Substituting into \eqref{eq:ccm},
\[
\rho_{AB} = \frac{(1+\rho)/4 - 1/4}{1/4} = \rho. \qedhere
\]
\end{proof}

So every value on the correlation scale is available with no classification problem between $A$
and $B$ at all. The examples that follow are more concrete, and each contradicts a specific
sentence in the paper.

\section{Eight examples}\label{sec:examples}

Each example needs only the inclusion rates $p_i$ and the co-occurrence rates $q_{ij}$, substituted
into \eqref{eq:ccm}.

Examples~\ref{ex:perfect} through \ref{ex:move} prescribe a law of prediction sets directly. By
Proposition~\ref{prop:law} the law of the sets is all the matrix depends on, so a law is all an
example needs to specify.

In those examples coverage is one, meaning that every set contains the true class. Any marginal
coverage lower bound is then met. We do not claim that randomized APS produces these laws. It
cannot produce coverage one, since with distinct scores its coverage is at most
$1 - \alpha + 1/(m+1)$.

Therefore Examples~\ref{ex:perfect} through \ref{ex:move} refute the authors' readings of the
matrix as claims about laws of prediction sets that meet the coverage bound. A claim restricted to
sets produced by APS needs a counterexample built inside APS. Example~\ref{ex:aps} is that
counterexample, which is why it comes first. The remaining examples are ordered from the simplest
to refute and the most absurd result downward.

In Examples~\ref{ex:rare} and \ref{ex:pool} a label $T$ is in every set. Its own entries are
undefined under the convention of Section~\ref{sec:statistic}, and the claims concern
$\rho_{AB}$ only.

\begin{example}[Opposite confusions, the same matrix, within randomized APS]\label{ex:aps}
Three classes with equal frequency. Two classifiers assign the probabilities $(0.6, 0.3, 0.1)$ to
the classes in the orders shown, so the true class is always ranked second.

\begin{center}
\begin{tabular}{lll}
\toprule
True class & Model 1 order & Model 2 order \\
\midrule
$A$ & $C > A > B$ & $B > A > C$ \\
$B$ & $A > B > C$ & $C > B > A$ \\
$C$ & $B > C > A$ & $A > C > B$ \\
\bottomrule
\end{tabular}
\end{center}

The two models get the same threshold. By \eqref{eq:aps} the score of the true label is
$0.9 - 0.3u$ under both models, uniform on $(0.6, 0.9)$. Calibrate both on the same inputs with the
same uniforms and both have the same threshold $\tau \in (0.6, 0.9)$. Write $r = (\tau - 0.6)/0.3$.

The two models produce sets with the same law. At a test input the top ranked class has score
$0.6 - 0.6U \le \tau$ and is always included. The second has score $0.9 - 0.3U$, which is at most
$\tau$ exactly when $U \ge (0.9 - \tau)/0.3$, an event of probability $r$. The third has score
$1 - 0.1U > 0.9 > \tau$ and is never included. Under either model, therefore, each singleton
$\set{A}$, $\set{B}$, $\set{C}$ occurs with probability $(1-r)/3$ and each pair $\set{A, B}$,
$\set{B, C}$, $\set{C, A}$ with probability $r/3$. The same law of sets gives the same matrix, and
the same coverage $r$, whose average over calibration draws is at least $1 - \alpha$.

The two models confuse classes in opposite directions. Model 1 labels every $A$ as $C$, every $B$
as $A$ and every $C$ as $B$. Model 2 runs the cycle the other way. Their confusion matrices are
transposes of each other, both with accuracy zero.

The authors write that the matrix ``contains a superset of the information conveyed by the confusion
matrix''. Here the sets are produced by randomized APS exactly as specified, the matrix is the same
for both models, and the confusion matrices differ. The matrix does not contain the confusion
matrix.
\end{example}

\begin{example}[A perfect classifier and a perfectly wrong one share a matrix]\label{ex:perfect}
This example makes the same point without the procedure, and with the accuracy gap as wide as it
can be. Three classes $A$, $B$, $C$ with equal frequency, and sets that always contain the true
class and one neighbor.

\begin{center}
\begin{tabular}{llcc}
\toprule
True class & Prediction set & Good model ranks first & Bad model ranks first \\
\midrule
$A$ & $\set{A, B}$ & $A$ & $B$ \\
$B$ & $\set{B, C}$ & $B$ & $C$ \\
$C$ & $\set{C, A}$ & $C$ & $A$ \\
\bottomrule
\end{tabular}
\end{center}

Each label is in two of the three sets, so every $p_i = 2/3$ and every variance is $2/9$. Each pair
shares exactly one of the three sets, so every $q_{ij} = 1/3$. Every off-diagonal entry is
\[
\rho_{ij} = \frac{1/3 - (2/3)(2/3)}{2/9} = -\tfrac12 .
\]

The good model has accuracy one and the bad model has accuracy zero. The labels and the sets are
identical for both, so the matrix is identical for both. The confusion matrix
$\Prob(Y = a, \hat Y = b)$ of the good model has all its mass on the diagonal and that of the bad
model has none there. The matrix does not contain ``a superset of the information conveyed by the
confusion matrix''.
\end{example}

\begin{example}[Zero correlation with 99\% directional confusion]\label{ex:zero}
Three classes and a small number $\varepsilon$, say $0.01$. Inputs arrive in four kinds.

\begin{center}
\begin{tabular}{lllc}
\toprule
Probability & True class & Prediction set & Ranked first \\
\midrule
$(1-\varepsilon)/2$ & $A$ & $\set{A, B}$ & $B$ \\
$\varepsilon/2$     & $A$ & $\set{A}$    & $A$ \\
$(1-\varepsilon)/2$ & $B$ & $\set{B}$    & $B$ \\
$\varepsilon/2$     & $C$ & $\set{C}$    & $C$ \\
\bottomrule
\end{tabular}
\end{center}

$A$ is in the sets of the first two kinds, so $p_A = (1-\varepsilon)/2 + \varepsilon/2 = 1/2$. $B$ is
in the sets of the first and third kinds, so $p_B = 1 - \varepsilon$. They appear together only in
the first kind, so $q_{AB} = (1-\varepsilon)/2 = p_A p_B$. The numerator of \eqref{eq:ccm} is zero,
so
\[
\rho_{AB} = 0 .
\]

For Bernoulli indicators that are not constant, zero correlation is independence. So $z_A$ and
$z_B$ are independent here. Yet among inputs whose true class is $A$, the classifier ranks $B$ first
with probability $1-\varepsilon$, which is 99\% at $\varepsilon = 0.01$.

The authors' Remark~3 says that unrelated classes give an entry near zero. A reader of a heat map
uses the converse, that an entry near zero means the classes are unrelated. The converse fails.
Independence of set membership does not imply absence of classification confusion.
\end{example}

\begin{example}[Perfect correlation with one in a million co-occurrence]\label{ex:rare}
Every set contains the true class $T$, ranked first. Two other labels $A$ and $B$ are added to the set
together with probability $\varepsilon$, and otherwise neither is added. Then $z_A = z_B$ on every
input, so
\[
\rho_{AB} = 1 \qquad \text{for every } 0 < \varepsilon < 1 .
\]

The probability that $A$ and $B$ appear together in a set is $q_{AB} = \varepsilon$, which may be
$0.4$, or $0.01$, or one in a million. The authors write in their abstract that the matrix
``captures the likelihood of class pairs co-occurring''. The likelihood is $\varepsilon$. The entry
is one regardless.
\end{example}

\begin{example}[Every pair positive from difficulty alone]\label{ex:hard}
Three classes with equal frequency, and difficulty independent of the true class. Half the inputs
are hard and receive the set of all three classes. The other half are easy and receive the singleton
set of their true class. The true class is ranked first in both cases, so accuracy is one.

Each label is in every hard set and in a third of the easy sets, so
\[
p_i = \tfrac12 \cdot 1 + \tfrac12 \cdot \tfrac13 = \tfrac23, \qquad p_i(1-p_i) = \tfrac29 .
\]
Two labels appear together only in hard sets, so $q_{ij} = 1/2$ for every pair, and
\[
\rho_{ij} = \frac{1/2 - (2/3)(2/3)}{2/9} = \tfrac14 .
\]

Every pair is positively correlated and the classifier is perfect. The correlation has one cause,
whether the input was hard, and says nothing about any pair. The authors' Remark~1 reads a positive
entry as ``intrinsically similar classes''.
\end{example}

\begin{example}[Pooling two clients creates a strong correlation from nothing]\label{ex:pool}
Every set contains the true class $T$, ranked first. On client~1 the labels $A$ and $B$ are each
added to the set independently with probability $0.9$. On client~2 they are each added independently
with probability $0.1$. The clients contribute equally to the pooled data.

Within each client $z_A$ and $z_B$ are independent, so $\rho_{AB} = 0$ on each. Pooled,
$p_A = p_B = (0.9 + 0.1)/2 = 1/2$, each variance is $1/4$, and
\[
q_{AB} = \tfrac12 (0.9 \times 0.9) + \tfrac12 (0.1 \times 0.1) = 0.41, \qquad
\rho_{AB} = \frac{0.41 - 1/4}{1/4} = 0.64 .
\]

The sign can be reversed. Let client~1 add $A$ with probability $0.9$ and $B$ with probability
$0.1$, and let client~2 do the opposite. Then $q_{AB} = \tfrac12(0.09) + \tfrac12(0.09) = 0.09$ and
$\rho_{AB} = (0.09 - 0.25)/0.25 = -0.64$.

The authors compute matrices across clients that are non identically distributed by construction
and read the entries as properties of the classifier. A strong entry of either sign is produced
here by the mixture of clients alone, with no confusion anywhere.
\end{example}

\begin{example}[Two classes the model never confuses get different entries]\label{ex:rates}
Three classes, cat, dog and plane, with equal frequency. The classifier returns the same set for every
input of a given class.

\begin{center}
\begin{tabular}{llccc}
\toprule
True class & Prediction set & $z_{\text{cat}}$ & $z_{\text{dog}}$ & $z_{\text{plane}}$ \\
\midrule
cat   & $\set{\text{cat}, \text{dog}}$ & 1 & 1 & 0 \\
dog   & $\set{\text{dog}}$             & 0 & 1 & 0 \\
plane & $\set{\text{plane}}$           & 0 & 0 & 1 \\
\bottomrule
\end{tabular}
\end{center}

A third of the inputs are cats, so $p_{\text{cat}} = 1/3$, and likewise $p_{\text{dog}} = 2/3$ and
$p_{\text{plane}} = 1/3$. Each variance is $2/9$. Cat and dog appear together on the cat inputs, so
$q_{\text{cat},\text{dog}} = 1/3$. Plane never appears with either, so the other two co-occurrence
rates are zero. Substituting,
\begin{align*}
\rho_{\text{cat},\text{dog}} &= \frac{1/3 - (1/3)(2/3)}{2/9} = \tfrac12, \\
\rho_{\text{cat},\text{plane}} &= \frac{0 - (1/3)(1/3)}{2/9} = -\tfrac12, \\
\rho_{\text{dog},\text{plane}} &= \frac{0 - (2/3)(1/3)}{2/9} = -1 .
\end{align*}

The authors' Remark~2 reads this as dog and plane being strongly dissimilar and cat and plane half
as dissimilar. But the classifier treats plane identically with respect to both. It never puts plane
in a set with either. The two entries differ only because dog is in twice as many sets as cat.
Lemma~\ref{lem:floor} gives the general statement.
\end{example}

\begin{example}[Moving one label occurrence between rows moves an entry from one to almost zero]\label{ex:move}
There are $n$ test inputs, and two rare labels $A$ and $B$ each appear in exactly one set. If they
appear in the same set then $z_A = z_B$ and $\rho_{AB} = 1$. If instead the one occurrence of $B$ is
in a different set then $p_A = p_B = 1/n$, $q_{AB} = 0$, and
\[
\rho_{AB} = \frac{0 - 1/n^2}{(1/n)(1-1/n)} = -\frac{1}{n-1},
\]
which at $n = 1000$ is about $-0.001$.

The authors present heat maps of entries with no uncertainty attached. For rare labels, moving a
single label occurrence from one row to another moves an entry across its whole range.
\end{example}

\section{Two lemmas}\label{sec:lemmas}

Write $o_i = p_i/(1-p_i)$ for the odds of inclusion of class $i$.

\begin{lemma}[The floor, and the attainable range]\label{lem:floor}
If classes $i$ and $j$ never share a set, then
\begin{equation}\label{eq:floor}
\rho_{ij} = -\sqrt{o_i\, o_j} = -\sqrt{\frac{p_i\, p_j}{(1-p_i)(1-p_j)}} .
\end{equation}
With $H$ classes of equal inclusion rate and singleton sets this equals $-1/(H-1)$. More generally,
for fixed inclusion rates $p_i \le p_j$ the entry $\rho_{ij}$ ranges over
\[
\Big[\, -\sqrt{o_i o_j},\ \sqrt{o_i / o_j}\, \Big] \quad \text{if } p_i + p_j \le 1, \qquad
\Big[\, -1/\sqrt{o_i o_j},\ \sqrt{o_i / o_j}\, \Big] \quad \text{if } p_i + p_j > 1 .
\]
\end{lemma}

\begin{proof}
If $q_{ij} = 0$ then \eqref{eq:ccm} reads
\[
\rho_{ij} = \frac{-p_i p_j}{\sqrt{p_i(1-p_i)p_j(1-p_j)}}
          = -\sqrt{\frac{p_i^2 p_j^2}{p_i(1-p_i)p_j(1-p_j)}}
          = -\sqrt{\frac{p_i p_j}{(1-p_i)(1-p_j)}} .
\]
With singleton sets every pair has $q_{ij} = 0$, and with equal inclusion rates $p_i = 1/H$, so
$o_i = 1/(H-1)$ and $\rho_{ij} = -1/(H-1)$.

For the range, $q_{ij}$ is the probability of the intersection of two events with probabilities $p_i
\le p_j$, so by the Fr\'echet bounds $\max(0, p_i + p_j - 1) \le q_{ij} \le p_i$. At the upper bound
the numerator of \eqref{eq:ccm} is $p_i - p_i p_j = p_i(1-p_j)$ and
\[
\rho_{ij} = \frac{p_i(1-p_j)}{\sqrt{p_i(1-p_i)p_j(1-p_j)}} = \sqrt{\frac{p_i(1-p_j)}{(1-p_i)p_j}} = \sqrt{o_i/o_j} .
\]
At the lower bound, if $p_i + p_j \le 1$ then $q_{ij} = 0$ and \eqref{eq:floor} applies. If
$p_i + p_j > 1$ then $q_{ij} = p_i + p_j - 1$, the numerator is $p_i + p_j - 1 - p_i p_j =
-(1-p_i)(1-p_j)$, and
\[
\rho_{ij} = \frac{-(1-p_i)(1-p_j)}{\sqrt{p_i(1-p_i)p_j(1-p_j)}} = -\sqrt{\frac{(1-p_i)(1-p_j)}{p_i p_j}} = -1/\sqrt{o_i o_j} . \qedhere
\]
\end{proof}

The lemma has two consequences for reading a heat map. The first concerns the scale. An entry
reads $+1$ only when $p_i = p_j$, so a rare class and a common class have a lower ceiling than two
common classes, and entries for pairs with different inclusion rates are not on a common scale.

The second concerns $\alpha$. With the model and calibration data fixed, the APS sets are nested in
$\alpha$. As $\alpha$ grows no inclusion rate increases, so every floor entry \eqref{eq:floor} moves
toward zero or stays put, whether or not the classifier changed, as long as the indicators involved
remain nondegenerate. The authors sweep $\alpha$ over four values and present the dependence as a
feature.

\begin{lemma}[Identification of the inclusion rates]\label{lem:ident}
Suppose every prediction set is a singleton and at least three classes have inclusion rate strictly
between zero and one. Then for any three distinct such classes $i$, $j$, $k$,
\begin{equation}\label{eq:ident}
o_i = -\frac{\rho_{ij}\,\rho_{ik}}{\rho_{jk}}, \qquad p_i = \frac{o_i}{1 + o_i} .
\end{equation}
The off-diagonal entries of the matrix therefore determine every inclusion rate exactly.
\end{lemma}

\begin{proof}
With singleton sets no two classes ever share a set, so \eqref{eq:floor} holds for every pair, and
$\rho_{ij}\rho_{ik} = \sqrt{o_i o_j}\sqrt{o_i o_k} = o_i \sqrt{o_j o_k} = -o_i \rho_{jk}$. Divide by
$-\rho_{jk}$, which is nonzero since $o_j, o_k > 0$. The map $o \mapsto o/(1+o)$ inverts
$p \mapsto p/(1-p)$.
\end{proof}

The authors argue that the matrix is privacy preserving because it ``only contains probabilities''
and so ``does not leak information on the size and composition'' of a client's data. Correlations
can be negative, so the matrix does not contain probabilities. A row normalized confusion matrix
does contain only probabilities, and reveals no counts either.

Lemma~\ref{lem:ident} says more. When every set is a singleton, the off-diagonal entries recover
every inclusion rate. When every set is the singleton of the ranked prediction, those rates are the
predicted class frequencies, which for an accurate classifier approximate the label histogram.

The matrix is offered as useful because it exposes how a client's data differs from the others. No
threat model, mechanism or leakage bound is given.

\section{What the coverage guarantee constrains}\label{sec:coverage}

The authors state that ``for a given input sample'' the set ``is guaranteed to contain the correct
class'' with probability close to $1-\alpha$. The guarantee \eqref{eq:coverage} is marginal. It
does not imply coverage conditional on the input, on a class or on a client. For continuous inputs
no nontrivial distribution free procedure delivers coverage conditional on every input
\citep{barber2021}.

The distinction matters here because the authors interpret individual class pairs and individual
federated clients as if validity transferred to them. It does not. Take ten equally weighted
clients, nine with coverage one and one with coverage zero. Pooled coverage is exactly $0.9$, which
meets a nominal $1 - \alpha = 0.9$, and the tenth client has no protection at all. The same
arithmetic applies to a class with a tenth of the data that is missed every time.

The authors also write the upper coverage term as $1/(n-1)$. For the randomized procedure with
almost surely distinct scores the term is $1/(n+1)$ \citep{romano2020, lei2018}, so theirs is looser
than necessary and probably a typographical error.

The federated experiments compare local matrices computed against a global calibration set across
clients that are non identically distributed by construction. Under that mismatch even the marginal
guarantee need not hold. No coverage or set sizes are reported for any client.

\section{What to report instead}\label{sec:instead}

If the question is which classes the classifier confuses, the probabilities answer it directly. The
soft confusion matrix
\[
C_{ab} = \E\big[\, p_b(X) \,\big|\, Y = a \,\big]
\]
is the average probability assigned to class $b$ on inputs of true class $a$. It is directional
and threshold free, and it needs no calibration set. It needs labeled evaluation data, as any
confusion matrix does.

The matrix $C_{ab}$ is a conditional mean, so it summarizes rather than determines. A model with $p_B(X) = 0.4$
on every input of class $A$ and a model with $p_B(X)$ equal to $0.6$ or $0.2$ with equal
probability have the same row $(0.6, 0.4)$, and error rates from $A$ to $B$ of zero and one half.
The confusion matrix of the ranked prediction carries that difference and $C_{ab}$ does not, so the
two are complements.

If labels are unavailable, the mean pairwise product $\E[p_i(X) p_j(X)]$ is symmetric and label
free. If sets are the object, because a downstream consumer receives sets and nothing else, then
the co-occurrence rate $q_{ij}$ is the quantity the abstract describes, and the coverage confusion
matrix of \citet{stutz2022}, which conditions on the true label, is directional.

In every case the inclusion rates and the set sizes should be reported alongside, since by
Lemma~\ref{lem:floor} an entry of \eqref{eq:ccm} cannot be read without them.

The authors compare against none of these. Their only comparison is with the hard confusion matrix,
which discards the probabilities entirely, so the contrast they draw is soft versus hard rather
than conformal versus not.

The authors' client selection experiment does not bear on any of this, but we note three things.
It selects the clients with the highest entries for two critical classes, while the centralized
section reads a falling entry as the classifier learning to separate them. Its Figures~6 and~7 are
captioned as matrix entries while their axes are class accuracy. Its Table~1 reports overall
CIFAR-100 accuracy falling from $0.511$ to $0.493$ under a heading of no effect, without
repetitions, seeds or error bars.

Simulations of Example~\ref{ex:aps}, of Examples~\ref{ex:hard} and \ref{ex:rates}, and of the
direction blindness of Example~\ref{ex:perfect}, on a synthetic ten class classifier with adaptive
prediction sets and with $C_{ab}$ and $\E[p_i p_j]$ alongside as $\alpha$ varies, are available at
\url{https://conformalprediction.net/demos/32-conformal-correlation.html}.

\section{Conclusion}\label{sec:conclusion}

The conformal correlation matrix is the phi coefficient of a thresholded softmax. By
Proposition~\ref{prop:law} it sees only the law of the sets, and by Proposition~\ref{prop:any} the
marginal coverage lower bound leaves any single entry free among the laws that satisfy it.

The eight examples show that each reading the authors attach to an entry, positive for similar,
negative for dissimilar, zero for unrelated, fails on a three class problem that meets that bound,
and Example~\ref{ex:aps} shows the failure inside randomized APS itself. Lemma~\ref{lem:floor}
shows that for pairs that never share a set the entry is a function of the inclusion rates alone,
and Lemma~\ref{lem:ident} shows that when every set is a singleton those rates can be read off the
matrix exactly. None of this is repaired by the certificate, which concerns the true label and says
nothing about which other labels sit beside it.

\bibliographystyle{plainnat}
\begin{thebibliography}{9}

\bibitem[Barber et~al.(2021)]{barber2021}
R.~F. Barber, E.~J. Cand\`es, A. Ramdas, and R.~J. Tibshirani.
\newblock The limits of distribution-free conditional predictive inference.
\newblock \emph{Information and Inference}, 10(2):455--482, 2021.

\bibitem[Lei et~al.(2018)]{lei2018}
J. Lei, M. G'Sell, A. Rinaldo, R.~J. Tibshirani, and L. Wasserman.
\newblock Distribution-free predictive inference for regression.
\newblock \emph{Journal of the American Statistical Association}, 113(523):1094--1111, 2018.

\bibitem[Perlo et~al.(2026)]{perlo2026}
A. Perlo, C.~F. Chiasserini, G. De~Veciana, and F. Malandrino.
\newblock Characterizing the performance of classification models through conformal correlation
matrices.
\newblock \emph{Computer Communications}, 247:108398, 2026.

\bibitem[Romano et~al.(2020)]{romano2020}
Y. Romano, M. Sesia, and E.~J. Cand\`es.
\newblock Classification with valid and adaptive coverage.
\newblock In \emph{Advances in Neural Information Processing Systems}, volume~33, pages 3581--3591,
2020.

\bibitem[Stutz et~al.(2022)]{stutz2022}
D. Stutz, K. Dvijotham, A.~T. Cemgil, and A. Doucet.
\newblock Learning optimal conformal classifiers.
\newblock In \emph{International Conference on Learning Representations}, 2022.

\bibitem[Vovk et~al.(2005)]{vovk2005}
V. Vovk, A. Gammerman, and G. Shafer.
\newblock \emph{Algorithmic Learning in a Random World}.
\newblock Springer, 2005.

\bibitem[Yule(1912)]{yule1912}
G.~U. Yule.
\newblock On the methods of measuring association between two attributes.
\newblock \emph{Journal of the Royal Statistical Society}, 75(6):579--652, 1912.

\end{thebibliography}

\end{document}
